\section{The Analysis for Difficult Graphs}\label{sec:difficult-graphs}
\label{sec:criterion}

Recall from Section~\ref{subsec:overview-graph} that a graph is difficult if it
satisfies \Hyp{} and is \loctree{\rad}{\nontree}; these are the graphs that
survive every reduction of Section~\ref{sec:graph-analysis}. Since a tree
neighborhood allows no
contradiction, we cannot rely on the local argument of
Section~\ref{sec:graph-analysis}. Instead, in this section we bound the
right-hand side of the main identity for such graphs, and then we prove the
main theorem.

We need the two bounds of Section~\ref{subsec:overview-three-conditions}, an
upper bound on \(\Kdef\) and an exponential-moment bound on \(\Tdef\), which are
of different types. The bound on \(\Kdef\) holds for every outcome of the
deletion experiment and only uses that most neighborhoods are trees, while the
bound on \(\Tdef\) is a probability bound that comes from the positive excess.

The bound on \(\Tdef\) is the harder one. As explained in
Section~\ref{subsec:overview-three-conditions}, we have
\(\Tdef\ge\cyc{\perc{\Gin}{1/2}}\) (see Section~\ref{subsec:long-to-cycle}
for the precise coupling), so it suffices to show that
\(\perc{\Gin}{1/2}\) has linear cycle rank. A direct first-moment count does not
give this. Since \(\cyc{H}=|E(H)|-|V(H)|+\ncomp{H}\), we have
\[
  \E\,\cyc{\perc{\Gin}{1/2}}
  =\frac{\medge}{2}-\nvert+\E\,\ncomp{\perc{\Gin}{1/2}} .
\]
When the density \(\medge/\nvert\) is below \(2\), the term
\(\medge/2-\nvert\) is negative, so the numbers of edges and vertices alone give
no lower bound. Everything then depends on the number of components, which is
hard to control directly. For this reason we take an indirect route: we bound
the probability of long kept paths from below and from above
(Section~\ref{subsec:long-path}), and compare the two bounds
(Section~\ref{subsec:long-to-cycle}).

\subsection{Standing assumptions and constants}\label{subsec:standing}

In this section and in Appendices~\ref{app:high-functional}
and~\ref{app:low-functional}, the constants \(\degb<\infty\) and \(\excd>0\) are
fixed, and every constant produced below depends only on them; in particular,
none of them depends on \(\rad\) or \(\nontree\), which are chosen at the end
(see Section~\ref{subsec:overview-three-conditions}). Recall the standing
hypothesis:

\medskip
\noindent\textbf{\Hyp.} \(\Gin\) has no self-loops, every vertex satisfies
\(2\le\dg{\Gin}{v}\le\degb\), and \(\exc{\Gin}\ge\excd\nvert\).

\medskip
We also use Definition~\ref{def:loctree} with several radii, and the next lemma
says that a smaller radius is always allowed.

\begin{lemma}[Monotonicity]\label{lem:loctree-monotone}
If \(\Gin\) is \loctree{\rad}{\nontree} and \(1\le R'\le \rad\), then
\(\Gin\) is \loctree{R'}{\nontree}.
\end{lemma}

\begin{proof}
The ball \(\ball{R'}{v}\) is an induced subgraph of \(\ball{\rad}{v}\), and it
is connected, because a shortest path from \(v\) to any of its vertices stays
inside it. Since a connected subgraph of a tree is a tree, \(\ball{R'}{v}\) is
a tree whenever \(\ball{\rad}{v}\) is. Hence every vertex with a non-tree ball
of radius \(R'\) also has a non-tree ball of radius \(\rad\).
\end{proof}

\medskip
\noindent\textbf{Notation of this section.} Besides the symbols of
Section~\ref{subsec:notation}, this section and Appendices~\ref{app:high-functional}
and~\ref{app:low-functional} use the following notation: the graph
\(\Akept\) of all kept edges, the retention probability \(\keeppz\) and the constant
\(\plow\) of Lemma~\ref{lem:high-functional}, the error term \(\thr\) of
Lemma~\ref{lem:low-functional}, and the cycle-rank density \(\cycd\) of
Lemma~\ref{lem:cycle-density}. Symbols that appear only in an appendix are
introduced there.

\subsection{The upper bound on \texorpdfstring{\(\Kdef\)}{K}}

\begin{Kboundrestatement}
If the \(\nvert\)-vertex graph \(\Gin\) is \loctree{\rad}{\nontree}, then
every outcome of the deletion experiment satisfies
\[
  \Kdef\le\left(\nontree+\frac{1}{2\rad}\right)\nvert .
\]
\end{Kboundrestatement}

This lemma needs neither the degree bound nor the excess.

\begin{proof}
Each forest component that contains a deferred edge gives a cycle, namely the
deferred edge together with the tree path between its endpoints, and cycles
from different forest components are vertex-disjoint. We call a cycle short if
its length is at most \(2\rad\), and long otherwise.

Every vertex of a short cycle has the whole cycle inside its ball of radius
\(\rad\), so this ball is not a tree. Since the short cycles are
vertex-disjoint, Definition~\ref{def:loctree} allows at most \(\nontree\nvert\)
of them. A long cycle has more than \(2\rad\) vertices, so, again by
disjointness, there are at most \(\nvert/(2\rad)\) long cycles. Adding the two
bounds proves the lemma.
\end{proof}

\subsection{The long-path functional}\label{subsec:long-path}

The bound on \(\Tdef\) uses the quantity \(\Lpath{\plen}{H}{\keepp}\) of
Section~\ref{subsec:overview-three-conditions}. Recall that
\(\perc{H}{\keepp}\) keeps each edge of \(H\) independently with probability
\(\keepp\), and that
\begin{equation}\label{eq:long-path-functional}
  \Lpath{\plen}{H}{\keepp}
  =\Prx{\perc{H}{\keepp}\text{ contains a path of length }\plen
  \text{ starting from a uniformly random vertex}} .
\end{equation}
The vertex is chosen independently of the edges, and the probability is over
both the random vertex and the random edges. A \emph{path of length \(\plen\)}
always means a self-avoiding path, that is, \(\plen+1\) distinct vertices joined
by \(\plen\) edges; walks that repeat a vertex are not counted.

Positive excess makes this quantity large, while a small cycle rank makes it
small, and together the two bounds below show that the cycle rank must be
linearly large.

\begin{lemma}[Positive excess gives a lower bound for long paths]
\label{lem:high-functional}
There are constants \(\keeppz=\keeppz(\excd)\in(0,1)\),
\(\plow=\plow(\degb,\excd)>0\) and \(\nontreez=\nontreez(\degb,\excd)>0\) with the
following property. Let \(\plen\ge1\) be an integer, and let \(\Gin\) satisfy
\Hyp{} and be \loctree{\plen}{\nontreez}. Then
\begin{equation}\label{eq:high-functional}
  \Lpath{\plen}{\Gin}{\keeppz/2}\ \ge\ \plow .
\end{equation}
\end{lemma}

\begin{lemma}[Cycle rank gives an upper bound for long paths]
\label{lem:low-functional}
Let \(\keepp<1\). For every \(\thr>0\) there is an integer
\(\plen=\plen(\degb,\keepp,\thr)\) such that every finite graph \(H\) with
maximum degree at most \(\degb\) satisfies
\begin{equation}\label{eq:low-functional}
  \Lpath{\plen}{H}{\keepp}
  \ \le\
  \thr+2\ballsz{\plen}\,\frac{\cyc{H}}{|V(H)|} .
\end{equation}
\end{lemma}

Lemma~\ref{lem:high-functional} is proved in
Appendix~\ref{app:high-functional}, and Lemma~\ref{lem:low-functional} in
Appendix~\ref{app:low-functional}.

Two points are important here. First, \(\plow\) does not depend on \(\plen\),
so we may first fix \(\plow\) and then choose \(\plen\) in
Lemma~\ref{lem:low-functional} with \(\thr=\plow/2\). Second, the lower bound
uses the retention probability \(\keeppz/2\), while the upper bound is applied
with \(\keeppz\); as the next subsection shows, the factor \(1/2\) comes from
the deletion experiment.

\subsection{From long paths to cycle rank}\label{subsec:long-to-cycle}

Let \(\Ffor\) be the forest of the first stage of the deletion experiment. Now
flip also the coins of the deferred edges, and let \(\Akept\) be the graph
formed by \(\Ffor\) and the kept deferred edges. Each edge of \(\Gin\) is then
kept independently with probability \(1/2\), so \(\Akept\) has the same law as
\(\perc{\Gin}{1/2}\). Note that \(\Akept\) is built from the same coins as
\(\Ffor\) and \(\Tdef\), which is what makes the following bound hold for every
outcome.

Since both endpoints of every
deferred edge lie in the same component of \(\Ffor\), adding the kept deferred
edges does not join components, and \(\ncomp{\Akept}=\ncomp{\Ffor}\). Also,
since \(\Ffor\) is a forest on all \(\nvert\) vertices, \(\cyc{\Ffor}=0\),
and
\[
  \cyc{\Akept}=|E(\Akept)|-\nvert+\ncomp{\Akept}
  =\cyc{\Ffor}+\bigl|E(\Akept)\setminus E(\Ffor)\bigr|
  =\bigl|E(\Akept)\setminus E(\Ffor)\bigr| .
\]
Every edge of \(E(\Akept)\setminus E(\Ffor)\) is deferred, so every outcome
satisfies
\begin{equation}\label{eq:T-nu}
  \Tdef\ \ge\ \cyc{\Akept} .
\end{equation}

Next, we thin once more, keeping each edge of \(\Akept\) independently with
probability \(\keepp\). Given \(\Akept\), the probability that a uniformly
random vertex starts a path of length \(\plen\) in the result is
\(\Lpath{\plen}{\Akept}{\keepp}\). On the other hand, each edge of \(\Gin\)
survives both steps independently with probability \(\keepp/2\), so averaging
over \(\Akept\) gives
\begin{equation}\label{eq:two-stage-functional}
  \Ex{\Lpath{\plen}{\Akept}{\keepp}}
  =\Lpath{\plen}{\Gin}{\keepp/2} .
\end{equation}

\begin{lemma}[Positive excess survives half-percolation]\label{lem:cycle-density}
There are constants \(\cycd=\cycd(\degb,\excd)>0\), an integer
\(\plenz=\plenz(\degb,\excd)\), and \(\nontreez=\nontreez(\degb,\excd)>0\) with
the following property. If \(\Gin\) satisfies \Hyp{} and is
\loctree{\plenz}{\nontreez}, then
\[
  \E\,\cyc{\perc{\Gin}{1/2}}\ \ge\ \cycd\,\nvert .
\]
\end{lemma}

\begin{proof}
Let \(\keeppz,\plow,\nontreez\) be the constants of
Lemma~\ref{lem:high-functional}, which depend only on \(\degb,\excd\). Apply
Lemma~\ref{lem:low-functional} with \(\keepp:=\keeppz\) and \(\thr:=\plow/2\),
and let \(\plenz\) be the resulting path length, which also depends only on
\(\degb,\excd\).

Assume \(\Gin\) satisfies \Hyp{} and is \loctree{\plenz}{\nontreez}. Then
Lemma~\ref{lem:high-functional} applies with \(\plen:=\plenz\), and
\[
  \plow
  \ \overset{\eqref{eq:high-functional}}{\le}\
  \Lpath{\plenz}{\Gin}{\keeppz/2}
  \ \overset{\eqref{eq:two-stage-functional}}{=}\
  \Ex{\Lpath{\plenz}{\Akept}{\keeppz}}
  \ \overset{\eqref{eq:low-functional}}{\le}\
  \frac{\plow}{2}+2\ballsz{\plenz}\,\frac{\E\,\cyc{\Akept}}{\nvert} .
\]
In the last step we apply~\eqref{eq:low-functional} to \(H=\Akept\) for every
outcome; this is allowed because \(\Akept\) has all \(\nvert\) vertices of
\(\Gin\) and maximum degree at most \(\degb\).
Rearranging gives \(\E\,\cyc{\Akept}\ge\cycd\nvert\) with
\[
  \cycd:=\frac{\plow}{4\ballsz{\plenz}} ,
\]
which depends only on \(\degb,\excd\).
\end{proof}

\subsection{The exponential moment of \texorpdfstring{\(\Tdef\)}{T}}

\begin{Tboundrestatement}
There are an integer \(\plenz=\plenz(\degb,\excd)\),
\(\nontreez=\nontreez(\degb,\excd)>0\) and \(\rate=\rate(\degb,\excd)>0\) such
that the following holds. If \(\Gin\) satisfies \Hyp{} and is
\loctree{\plenz}{\nontreez}, then
\[
  \Ex{\left(\tfrac34\right)^{\Tdef}}\le 2e^{-\rate\nvert} .
\]
\end{Tboundrestatement}

We take \(\plenz\) and \(\nontreez\) from Lemma~\ref{lem:cycle-density}. To
turn the bound on the expectation into a tail bound, we use McDiarmid's
bounded-differences inequality in the following form: if \(f\) is a function of
\(\medge\) independent random variables, and changing one of them changes \(f\)
by at most \(1\), then
\[
  \Prx{f\le\E f-t}\le\exp\left(-\frac{2t^2}{\medge}\right)
  \qquad(t>0).
\]

\begin{proof}
The graph \(\Gin\) has \(\medge\le\degb\nvert/2\) edges, and
\(\Akept=\perc{\Gin}{1/2}\) is determined by \(\medge\) independent fair coins.
Changing one coin changes \(\cyc{\Akept}\) by at most \(1\). Indeed, suppose
the coin adds an edge: if the edge joins two vertices in the same component,
the cycle rank goes up by \(1\), and if it joins two components, the number of
edges goes up by \(1\) and the number of components goes down by \(1\), so the
cycle rank does not change.

Lemma~\ref{lem:cycle-density} gives \(\E\,\cyc{\Akept}\ge\cycd\nvert\). We apply
the inequality above with \(f:=\cyc{\Akept}\) and \(t:=\cycd\nvert/2\), and use
\(\medge\le\degb\nvert/2\):
\[
  \Prx{\cyc{\Akept}<\tfrac{\cycd}{2}\nvert}
  \le\exp\left(-\frac{2(\cycd\nvert/2)^2}{\medge}\right)
  \le\exp\left(-\frac{\cycd^{2}}{\degb}\nvert\right) .
\]
By~\eqref{eq:T-nu}, the same bound holds for
\(\Prx{\Tdef<\tfrac{\cycd}{2}\nvert}\).

Now we split according to whether \(\Tdef\) is small:
\[
  \Ex{\left(\tfrac34\right)^{\Tdef}}
  \le \left(\tfrac34\right)^{\cycd\nvert/2}+\Prx{\Tdef<\tfrac{\cycd}{2}\nvert}
  \le \exp\left(-\frac{\cycd}{2}\log\frac43\cdot\nvert\right)
     +\exp\left(-\frac{\cycd^{2}}{\degb}\nvert\right)
  \le 2e^{-\rate\nvert},
\]
where
\[
  \rate:=\min\left\{\frac{\cycd}{2}\log\frac43,\ \frac{\cycd^{2}}{\degb}\right\}>0
\]
depends only on \(\degb,\excd\).
\end{proof}

\subsection{Proof of Theorem~\texorpdfstring{\ref{thm:overview-three-conditions}}{1.12}}

\begin{threeconditionsrestatement}
There are \(\radz=\radz(\degb,\excd)\) and \(\nontreeo=\nontreeo(\degb,\excd)>0\)
such that the following holds. If \(\rad\ge\radz\), \(\nontree\le\nontreeo\), and
the \(\nvert\)-vertex graph \(\Gin\) satisfies \Hyp{} and is
\loctree{\rad}{\nontree}, then
\[
  \E\,\solcomp{\rform{\Gin}}
  =\Ex{2^{\Kdef}\left(\tfrac34\right)^{\Tdef}}
  \le 2e^{-\rate\nvert/2},
\]
where \(\rate=\rate(\degb,\excd)>0\) is the rate of
Lemma~\ref{lem:overview-T-bound}.
\end{threeconditionsrestatement}

\begin{proof}
The equality is the main identity, Theorem~\ref{thm:overview-main-identity},
which applies because \Hyp{} forbids self-loops.

Let \(\plenz,\nontreez,\rate\) be the constants of
Lemma~\ref{lem:overview-T-bound}; they depend only on \(\degb,\excd\). Set
\[
  \radz:=\max\left\{\plenz,\ \left\lceil\frac{2\log2}{\rate}\right\rceil\right\},
  \qquad
  \nontreeo:=\min\left\{\nontreez,\ \frac{\rate}{4\log2}\right\} .
\]
If \(\rad\ge\radz\) and \(\nontree\le\nontreeo\), then
\[
  \Bigl(\nontree+\frac1{2\rad}\Bigr)\log2
  \le\frac{\rate}{4}+\frac{\rate}{4}
  =\frac{\rate}{2} .
\]

Moreover, since \(\rad\ge\plenz\) and \(\nontree\le\nontreez\),
Lemma~\ref{lem:loctree-monotone} shows that \(\Gin\) is
\loctree{\plenz}{\nontreez}, so Lemma~\ref{lem:overview-T-bound} applies.
Since Lemma~\ref{lem:overview-K-bound} gives
\(\Kdef\le(\nontree+\frac1{2\rad})\nvert\) for every outcome, we get
\[
  \Ex{2^{\Kdef}\left(\tfrac34\right)^{\Tdef}}
  \le 2^{(\nontree+\frac1{2\rad})\nvert}\,\Ex{\left(\tfrac34\right)^{\Tdef}}
  \le 2\cdot e^{\left((\nontree+\frac1{2\rad})\log2-\rate\right)\nvert}
  \le 2e^{-\rate\nvert/2} .
\]
\end{proof}

\subsection{Proof of the main theorem}\label{subsec:main-proof}

All parts are now available, namely the reductions of
Section~\ref{sec:graph-analysis}, the identity of
Section~\ref{sec:solution-components}, and
Theorem~\ref{thm:overview-three-conditions}.

\begin{mainrestatement}
Fix \(\excpar>0\). Then every finite multigraph \(\Gin\) with \(\nvert\)
vertices and
\[
  \exc{\Gin}\ge\excpar\nvert
\]
satisfies
\[
  \Prx{\rform{\Gin}\SAT}\le e^{-\Omega(\nvert)},
\]
where the implied constant depends only on \(\excpar\). In particular the
satisfiability probability tends to \(0\) as \(\nvert\to\infty\).
\end{mainrestatement}

\begin{proof}[Proof of Theorem~\ref{thm:main}]
Let \(\degb\) and \(\excd\) be the constants of
Lemma~\ref{lem:overview-bounded-core}, which depend only on \(\excpar\), and
let \(\radz\) and \(\nontreeo\) be the constants of
Theorem~\ref{thm:overview-three-conditions} for these \(\degb,\excd\). In this
proof, \(H_0,H_1,H_2\) are three graphs, and every \(\Omega(\cdot)\) hides a
constant that depends only on \(\excpar\).

\emph{Step 1.} If \(\medge\ge2.41\,\nvert\), then
Lemma~\ref{lem:overview-first-moment} gives
\(\Prx{\rform{\Gin}\SAT}\le e^{-\Omega(\nvert)}\) and we are done. Otherwise
Lemma~\ref{lem:overview-bounded-core} gives a graph \(H_0\) with
\[
  |V(H_0)|\ge\frac{\excpar}{\degb}\,\nvert,
  \qquad
  2\le\dg{H_0}{v}\le\degb,
  \qquad
  \exc{H_0}\ge2\excd|V(H_0)|,
\]
and \(\Prx{\rform{\Gin}\SAT}\le\Prx{\rform{H_0}\SAT}\).

\emph{Step 2.} We apply Lemma~\ref{lem:overview-positive-components} to
\(H_0\) and call the result \(H_1\). Its degrees are the same as in \(H_0\), and
\[
  \exc{H_1}\ge\exc{H_0}\ge2\excd|V(H_0)|\ge2\excd|V(H_1)|,
  \qquad
  |V(H_1)|\ge\frac2\degb\exc{H_0}\ge\frac{4\excd}{\degb}|V(H_0)|=\Omega(\nvert) .
\]
Every component of \(H_1\) without a self-loop has cycle rank at least \(3\),
and every component with a self-loop has cycle rank at least \(2\).

\emph{Step 3.} We apply Lemma~\ref{lem:remove-loops} to \(H_1\). Either
\(\Prx{\rform{H_1}\SAT}\le e^{-\Omega(|V(H_1)|)}\), or we get a graph \(H_2\)
without self-loops with
\[
  |V(H_2)|\ge\frac{2\excd}{\degb}|V(H_1)|=\Omega(\nvert),
  \qquad
  2\le\dg{H_2}{v}\le\degb,
  \qquad
  \exc{H_2}\ge\excd|V(H_2)|,
\]
in which every component has cycle rank at least \(3\), and
\(\Prx{\rform{H_1}\SAT}\le\Prx{\rform{H_2}\SAT}\).

\emph{Step 4.} If \(H_2\) is not \loctree{\radz}{\nontreeo}, then
Lemma~\ref{lem:overview-finite-counting}, with \(\rad=\radz\) and
\(\nontree=\nontreeo\), gives a radius \(\radp\) and a fraction \(\pack\), and
Corollary~\ref{cor:overview-packing} gives
\(\Prx{\rform{H_2}\SAT}\le e^{-\Omega(|V(H_2)|)}\). Otherwise \(H_2\) satisfies
\Hyp{} and is \loctree{\radz}{\nontreeo}, so
Corollary~\ref{cor:overview-sat-bound} and
Theorem~\ref{thm:overview-three-conditions}, with \(\rad=\radz\) and
\(\nontree=\nontreeo\), give
\[
  \Prx{\rform{H_2}\SAT}
  \le \E\,\solcomp{\rform{H_2}}
  \le e^{-\Omega(|V(H_2)|)} .
\]

In every case the bound is exponential in \(|V(H_1)|\) or \(|V(H_2)|\), both of
which are \(\Omega(\nvert)\), and all implied constants depend only on
\(\excpar\). Together with the chain
\(\Prx{\rform{\Gin}\SAT}\le\Prx{\rform{H_0}\SAT}\le\Prx{\rform{H_1}\SAT}
\le\Prx{\rform{H_2}\SAT}\), this gives
\[
  \Prx{\rform{\Gin}\SAT}\le e^{-\Omega(\nvert)} .
\]
\end{proof}
